\chapter{Introduction}\label{ch01:chap}

This monograph develops, in full and at reference depth, a deliberately small
executable kernel for the typed-composition layer of the Tootra Knowledge
System (TKS) and its secular derivative, the Cognitive Architecture System
(CAS). The kernel --- TKS v8.0 --- supplies named finite objects with
explicitly ordered carriers, total finite-function arrows with declared source
and target, partial type-checked categorical composition, a hard type
separation between one-carrier endomaps and heterogeneous semigroupoid paths, a
deterministic standard-basis lift into finite-dimensional Hilbert spaces, and
stable denotational normal forms computed by shortlex minimization over a
completed finite closure. Everything the kernel does is either proved here,
tested by bounded exhaustive enumeration, kernel-checked in a pinned proof
assistant, or explicitly recorded as open. Nothing in this monograph is
peer-reviewed; the artifact status of the entire program is \textsc{draft /
gate 1 prototype / ring 4 / not licensed / not frozen}, and we will repeat that
status wherever a reader might otherwise infer more.

The present chapter explains why so small a kernel deserves so long a book.
The short answer is that the kernel exists to make three specific reasoning
failures \emph{impossible to write down}, and that the discipline required to
achieve this --- in the mathematics, in the implementation, and in the evidence
--- turns out to be instructive far beyond the kernel's size. The three
failures are untyped composition, invalid power inference, and nondeterministic
identity. Each is developed below with an extended worked scenario showing
exactly how the failure produces a wrong conclusion from locally reasonable
steps. Each scenario is then revisited at the end of the chapter to show which
kernel rule makes the failing step a type error rather than a silent
miscomputation.

\section{The TKS/CAS formal program and its compiler}\label{ch01:sec:program}

\subsection{The ontology in brief}

The TKS/CAS formal program, in its v7.4 form and the v7.5 through v7.7
extensions, defines a rich ontology for describing cognitive architecture: ten
Noetic operators, four processing Domains, seven Core Drives, four
Registers introduced in the later extensions, and a disciplined operation
notation in which compound cognitive descriptions are written as expressions
over these primitives. The ontology is served by a production notation
compiler that parses, validates, and evaluates expressions of the notation.
This monograph takes the ontology entirely as \emph{context}: the kernel
developed here deliberately encodes none of it. The refusal is not modesty
for its own sake; it is a gate. Mapping the canonical Noetics or any CAS
construct onto kernel arrows is a separately tracked research obligation
(PO-TKS8-CORE-08 in the ledger of Appendix~B), it has not been performed, and
it is recorded as \textsc{open / not licensed}. Until that mapping is
reviewed, every statement in this book is a statement about named finite sets
and total functions --- no more.

\subsection{Kinds in the production compiler}

The production notation compiler enforces what we will call \emph{value
kinds}. Each value flowing through an expression carries a tag drawn from a
small fixed vocabulary --- state, drive, relation, program, and so on --- and
the compiler rejects expressions that combine kinds in ways the grammar
forbids. Kind checking is genuinely useful: it stops a drive value from being
used where a relation is expected, and it catches a large class of
transcription errors at parse time.

What a kind system does not do --- cannot do, at this granularity --- is track
\emph{categorical type}: the exact source object and exact target object of
each operation. Two operations may both carry the kind ``relation'' while one
of them is a map out of a seven-element carrier and the other is a map into a
four-element carrier. The kind system sees two relations and permits the
combination; the composition is nevertheless meaningless, because the
intermediate objects disagree. In categorical language
\cite{maclanecwm}: a kind system checks membership in a class of arrows,
whereas composition legality is a statement about the \emph{objects} at the
interface, and objects are precisely what kinds erase.

\subsection{The artifact and its evidence infrastructure}

Because later chapters constantly refer to the kernel's evidence
apparatus, we fix the vocabulary here. The \emph{kernel} is the Python
package \texttt{src/tks\_v8\_core}, comprising three modules:
\texttt{category.py} (objects, arrows, composition, endomaps, paths),
\texttt{deterministic\_lift.py} (the standard-basis lift and its derived
operators), and \texttt{normal\_forms.py} (generator catalogs, typed
words, the shortlex closure). A \emph{manifest} binds the exact bytes of
these modules and their test suites by SHA-256 hash. The \emph{Gate~1
verifier} (\texttt{scripts/verify\_tks\_v8\_gate1.py}) checks artifact
hashes against the manifest and reruns the executable suite; the
\emph{Lean verifier} (\texttt{scripts/verify\_psr8\_lean\_gate1.py})
rebuilds a pinned, isolated Lean~4.32.1 formal seed \cite{lean4} under an
audited axiom footprint. ``Gate~1'' names the project's first review
gate: the kernel is a Gate~1 \emph{candidate} --- the manifest does not
hash itself, no reviewed commit or acceptance receipt yet binds the
manifest hash, and until one does, nothing here is a freeze receipt.
The validation snapshot of 2026-08-04 records \textsc{pass} on all 20
tests of the suite.

Two terms of art recur. A claim's \emph{evidential status} is its entry
in the proof-obligation ledger (Appendix~B), drawn from a fixed
vocabulary explained in Remark~\ref{ch01:rem:status} below. The
\emph{ring/license} annotations (\textsc{ring 4 / not licensed}) are the
program's admission-control labels: they mark the entire artifact as an
unreviewed research prototype whose claims may not be transferred into
any production, marketing, or canonical-ontology context. This
monograph inherits both disciplines wholesale and adds no claim beyond
what the ledger registers.

\subsection{Provenance of the failure analysis}

The three failure modes below were identified by the Project Coherence
mathematical red-team audit of the TKS/CAS formal layer. The audit's findings
are inherited here as the kernel's requirements; we restate them in our own
notation and expand each into a worked scenario. The scenarios are
constructed for this monograph --- they are illustrative models of the failure
classes, not transcripts of observed production incidents, and no claim about
the behavior of any deployed system is made or implied.

\begin{finding}[Audit finding, restated]\label{ch01:finding:audit}
The production compiler's kind system does not enforce categorical source and
target objects. Consequently three classes of invalid inference are
expressible without error: (i) composition of operations whose intermediate
objects disagree; (ii) application of single-operator iteration theorems to
heterogeneous generator sequences; (iii) dependence of canonical
representatives and digests on interpreter hash state. TKS v8.0 is the
response: the smallest kernel in which each class is a type error.
\end{finding}

\section{Kinds are not types}\label{ch01:sec:kinds}

Before the scenarios, it is worth being precise about the gap, because the gap
is the entire motivation for the kernel.

A \emph{kind}, as the production compiler uses the term, is a unary predicate
on values: ``this value is a relation.'' A \emph{categorical type}, as the
kernel uses the term, is a pair of objects: ``this arrow goes from $X$ to
$Y$,'' written $f : X \to Y$. The difference matters at exactly one place:
composition. Kind-level checking validates $g \circ f$ whenever $f$ and $g$
have composable kinds. Type-level checking validates $g \circ f$ only when
the declared target object of $f$ \emph{equals} the declared source object of
$g$ --- equality of objects, not of kinds, not of cardinalities, not of
element sets.

Three progressively stricter equality tests could be imagined at the
interface, and the kernel adopts the strictest:

\begin{enumerate}
\item \emph{Cardinality equality}: $|Y| = |Y'|$. This licenses the index
coincidence failure of Scenario~I below and is rejected.
\item \emph{Carrier equality}: the carriers of $Y$ and $Y'$ are equal as
ordered sequences. This still coerces semantically distinct objects that
happen to share elements --- two registers both carrying $(0,1)$ --- and is
rejected.
\item \emph{Name-and-carrier equality}: the declared names and the ordered
carriers both agree. This is the kernel's rule
(Chapter~\ref{ch03:chap}), and the anti-coercion argument for it is given at
length in Section~\ref{ch03:sec:name-identity}.
\end{enumerate}

A natural question is why the response to the gap is a separate kernel
rather than an extension of the production compiler's kind system into a
full type system. The answer is procedural as much as technical. The
production compiler is a live artifact with existing users and existing
semantics; retrofitting object-level types into it would change the
acceptance behavior of existing expressions, and doing so \emph{before}
the categorical layer has been independently specified, proved, and
reviewed would repeat the exact mistake the audit flagged --- semantics
adopted by implementation accident rather than by reviewed decision. The
kernel therefore stands beside the compiler, not inside it: a reference
model whose small size makes complete specification and bounded
exhaustive validation feasible, and whose integration with the compiler
is deliberately held open as a separately gated obligation
(PO-TKS8-CORE-09, \textsc{open}). If and when integration is designed,
the kernel provides the semantics to integrate \emph{to}; until then it
provides something almost as useful --- a precise statement of what the
kind system does not check, with executable witnesses.

The remainder of this chapter shows what goes wrong when the interface check
is weaker than test (3), and, in the third scenario, what goes wrong when even
test (3) is undermined by nondeterministic carrier order.

\section{Failure mode I: untyped composition}\label{ch01:sec:untyped}

\subsection{The failure class}

Nothing in a kind system prevents composing operations whose intermediate
objects disagree. The result is a well-formed expression --- it parses, it
kind-checks, it may even execute --- with no denotation. The most dangerous
member of this class is not the case where execution crashes; it is the case
where a cardinality coincidence lets an index-based implementation run to
completion and return a value. The value is then read as if it meant
something, and a conclusion is drawn that no step of the computation ever
licensed.

\subsection{Worked scenario}\label{ch01:sec:untyped-scenario}

We construct the scenario with kernel-style objects so that the reader can
see exactly which datum the weak checker ignores. Take three finite objects:

\begin{itemize}
\item $D = (\texttt{DriveProfile}, (d_1, d_2, d_3, d_4, d_5, d_6, d_7))$, a
seven-element carrier whose positions are read, in the surrounding analysis,
as seven drive labels;
\item $E = (\texttt{EvalSummary}, (e_0, e_1, e_2, e_3))$, a four-element
carrier of evaluative summary labels;
\item $Q = (\texttt{QuadState}, (q_0, q_1, q_2, q_3))$, a four-element
carrier of unrelated quadrant labels used by a different analysis pipeline.
\end{itemize}

Two total maps are on hand, both carrying the kind ``relation'' in a
kind-checked setting:

\begin{itemize}
\item $p : D \to E$, a summarization step; say $p$ has mapping tuple
$(0, 0, 2, 2, 2, 1, 3)$, meaning $p(d_1) = e_0$, $p(d_2) = e_0$,
$p(d_3) = e_2$, and so on.
\item $r : Q \to D$, a reading step from the quadrant pipeline; say $r$ has
mapping tuple $(4, 6, 5, 0)$, meaning $r(q_0) = d_5$, $r(q_1) = d_7$,
$r(q_2) = d_6$, $r(q_3) = d_1$.
\end{itemize}

An analyst wants to chase a drive profile through summarization and back into
drive space, and writes the composite $r \circ p$. The composite is ill-typed:
the target of $p$ is $E$ and the source of $r$ is $Q$, and $E \neq Q$ ---
different names, different carriers. But both carriers have four elements. An
implementation that represents maps as bare index tuples, with no declared
source and target, evaluates $r \circ p$ without complaint: it takes the
output index of $p$ and uses it to index into the table of $r$.

Run the numbers. Start from $d_3$. Then $p(d_3) = e_2$, held internally as
index $2$. The untyped composition feeds index $2$ into the table of $r$,
whose entry at position $2$ is $5$ --- but position $2$ of the table of $r$ is
the row for $q_2$, not for $e_2$. The machine returns $d_6$. The analyst
reads off: ``chasing $d_3$ through summarization and reading yields $d_6$,''
and proceeds to build on the association between $d_3$ and $d_6$.

The conclusion is wrong in the strongest possible sense: not incorrect by a
computable error term, but \emph{underived}. No step of any analysis ever
asserted what the reading map $r$ does to the evaluative label $e_2$, because
$r$ is not defined on evaluative labels at all. The number $5$ that surfaced
mid-computation is the image of $q_2$, an element of a carrier from an
unrelated pipeline, smuggled in through nothing but the coincidence
$|E| = |Q| = 4$. Had $|E|$ been $5$, the composition would have crashed on
some inputs and the error would have been caught; the four-versus-four
coincidence is exactly what makes the failure silent.

Note carefully what each checking regime says about this composite:

\begin{itemize}
\item \emph{Kind checking}: $p$ and $r$ are both relations; relations compose;
accepted.
\item \emph{Cardinality checking}: $|E| = |Q| = 4$; accepted.
\item \emph{Carrier checking}: $(e_0, e_1, e_2, e_3) \neq (q_0, q_1, q_2,
q_3)$; rejected --- but only because the element sets happen to differ. Rename
the quadrant labels to $e_0, \dots, e_3$ during a refactor and carrier
checking silently accepts again.
\item \emph{Kernel checking}: the target object of $p$ is
$(\texttt{EvalSummary}, \dots)$ and the source object of $r$ is
$(\texttt{QuadState}, \dots)$; the names differ; the composition raises a
typed error naming both objects. No refactor of element spellings can
reintroduce the silent acceptance without also merging the declared object
identities, which is a visible, reviewable act.
\end{itemize}

\subsection{A taxonomy of the failure's presentations}

It is worth separating the three ways an untyped composition can present,
because they have very different detection profiles:

\begin{enumerate}
\item \emph{Loud}: the cardinalities disagree and some input's index
falls outside the second table. The composition crashes --- on those
inputs, in that run. Detection is likely but not certain: if the
offending inputs never occur in testing, the crash ships.
\item \emph{Silent-total}: the cardinalities agree, as in the worked
scenario. Every input executes, every output is a well-formed element,
and no observable signal distinguishes the underived result from a
correct one. Detection requires semantic review of a computation that
looks healthy.
\item \emph{Silent-partial}: the cardinalities disagree but the
observed inputs happen to land inside the smaller table. The pipeline
runs clean for weeks and crashes on the first unusual input ---
combining the worst of both: a latency-bombed crash \emph{and} a
history of silently underived outputs before it.
\end{enumerate}

Only (1) is reliably caught by ordinary testing, and an adversarially
unlucky data distribution converts (1) into (3). A type check at
composition time collapses all three into a fourth presentation:
\emph{immediate, deterministic, input-independent rejection at
construction} --- the only presentation whose detection probability is
$1$ and whose detection latency is zero.

\subsection{The lesson}

The failure is not exotic. Index-tuple representations of finite maps are the
natural implementation choice, and cardinality coincidences among small
carriers are the norm, not the exception --- the TKS/CAS ontology alone has a
four-element Domain axis, a four-element Register axis, and numerous
four-valued summaries. Any two of these will silently interchange under an
untyped composition. The kernel's response (Chapter~\ref{ch03:chap} and
Chapter~4) is that an arrow \emph{is} its judgment: source object, target
object, and mapping tuple travel together, and composition is a partial
operation defined only when the middle objects are equal as named, ordered
objects.

\section{Failure mode II: invalid power inference}\label{ch01:sec:power}

\subsection{The failure class}

Finite endomaps obey strong iteration theorems. The one this monograph
proves as Theorem~5.1 (Chapter~5; the underlying mathematics is
Theorem~\ref{ch02:thm:factorial-idempotent} of the primer) is the safe
finite power: for any total $f : X \to X$ on a finite carrier, the power
$f^{|X|!}$ is idempotent, $f^{|X|!} \circ f^{|X|!} = f^{|X|!}$. Results of
this shape license fixed-point and stabilization conclusions: iterate long
enough and the process reaches an idempotent regime.

Every such theorem has a hypothesis that informal operator notation makes
easy to lose: the iterated thing must be \emph{one} endomap --- a single
element of the finite monoid $\mathrm{End}(X)$, applied repeatedly. A
workflow that applies generators $g_1, g_2, g_3, \dots$ in temporal order is
not, in general, iterating one endomap. If the $g_t$ differ from step to
step, the composite $g_t \circ \cdots \circ g_1$ is an element of a
\emph{path} through the category, and if the state spaces change over time it
is not even an endomap. Writing the sequence as ``$g^t$'' is a notational
pun, and importing an idempotence theorem through that pun produces
conclusions that are false --- not approximately false, but exactly false,
refutable on a three-element carrier.

\subsection{Worked scenario}\label{ch01:sec:power-scenario}

Let $X$ be a three-element carrier, positions $\{0, 1, 2\}$. A session
applies two alternating generators, both honest endomaps of $X$:

\begin{itemize}
\item $a$, the cyclic rotation with mapping tuple $(1, 2, 0)$: it sends
$0 \mapsto 1$, $1 \mapsto 2$, $2 \mapsto 0$;
\item $b$, the swap with mapping tuple $(1, 0, 2)$: it exchanges $0$ and
$1$ and fixes $2$.
\end{itemize}

The session trace is $a, b, a, b, a, b, \dots$: at odd steps apply $a$, at
even steps apply $b$. Write $s_t$ for the composite state map after $t$
steps, so $s_1 = a$, $s_2 = b \circ a$, $s_3 = a \circ b \circ a$, and in
general $s_t$ is the left-to-right alternating product.

An analyst models the session as ``the session operator applied $t$ times''
--- a single $f$ with $s_t = f^t$ --- and invokes the safe finite power with
$|X| = 3$, hence $|X|! = 6$: $f^6$ is idempotent, so
$f^{12} = f^6 \circ f^6 = f^6$, and the analyst concludes that the state map
after twelve steps equals the state map after six steps. \emph{The process
has stabilized by step six}; monitoring beyond that point is declared
redundant.

Now compute what actually happens. The two-step round is
\[
h = b \circ a : \quad
0 \xrightarrow{\;a\;} 1 \xrightarrow{\;b\;} 0, \qquad
1 \xrightarrow{\;a\;} 2 \xrightarrow{\;b\;} 2, \qquad
2 \xrightarrow{\;a\;} 0 \xrightarrow{\;b\;} 1,
\]
so $h$ has mapping tuple $(0, 2, 1)$: it fixes $0$ and swaps $1$ and $2$.
Then $s_6 = h^3$ and $s_{12} = h^6$. Since $h$ is a transposition,
$h^2 = \id_X$, hence
\[
s_6 = h^3 = h \neq \id_X = h^6 = s_{12}.
\]
The predicted equality $s_{12} = s_6$ is false: after six steps the state
map swaps positions $1$ and $2$; after twelve it is the identity. Any
downstream conclusion keyed to ``the stabilized map'' --- for instance, that
position $1$ has permanently migrated to position $2$ --- is refuted six
steps later by the process itself.

The instructive point is that \emph{every hypothesis of the theorem holds for
each generator separately}. Both $a$ and $b$ are single finite endomaps;
$a^6 = \id_X$ and $b^6 = \id_X$, and the identity is idempotent, so the safe
finite power theorem is true of $a$ and true of $b$. It fails for the session
because the session is not a power of any single endomap: there is no
$f : X \to X$ with $f^t = s_t$ for all $t$ (already $s_1 = a$ and
$s_2 = b \circ a$ would force $f = a$ and $f^2 = b\circ a$, but
$a^2 = (2,0,1) \neq (0,2,1)$). The pun ``$g^t$'' asserted the existence of
such an $f$, silently, in the notation itself.

The heterogeneous variant is even sharper. If the generators cross carriers
--- $g_1 : X_0 \to X_1$, $g_2 : X_1 \to X_2$ with $X_0 \neq X_2$ --- then
``the session operator squared'' is not merely a false claim but an ill-typed
expression: the composite has type $X_0 \to X_2$, and composing it with
itself requires $X_2 = X_0$. There is no monoid in which the power lives.

\subsection{The lesson}

The kernel's response (Chapters~5 and~6) is a type-level separation: the
\texttt{Endomap} constructor is available only for arrows whose source and
target are the \emph{same} object, and only \texttt{Endomap} carries a power
operation. Heterogeneous or time-varying sequences are represented as
\texttt{SemigroupoidPath} values, which have a typed composite and
\emph{deliberately no power operation} --- the method is absent, and the
runtime types are distinct, so a path cannot be passed where an endomap is
required. The invalid inference of this scenario is not discouraged by
documentation; it is unwritable. One cannot ask the kernel for $s^{6}$
because the object $s$ on which one would ask it does not have the method,
and constructing an \texttt{Endomap} from the session requires exhibiting the
single arrow that the pun only pretended exists.

\section{Failure mode III: nondeterministic identity}\label{ch01:sec:nondet}

\subsection{The failure class}

The first two failures are semantic. The third is evidential, and it
undermines even a correctly typed system if left unaddressed. Everything in
the kernel downstream of an object --- basis indices of the Hilbert-space
lift, mapping tuples of arrows, shortlex ranks of words, canonical
representatives of denotations, and every digest computed over any of these
--- depends on the \emph{order} of the object's carrier. If that order is
left to the iteration order of a hash-based container, it depends on
interpreter state: in CPython, on the per-process string hash randomization
controlled by \texttt{PYTHONHASHSEED}. Two runs of the same program on the
same input may then disagree about which element sits at index $2$, and
therefore about every downstream artifact, and therefore about every
downstream digest. Replayability --- the ability of an auditor to rerun a
computation and obtain bit-identical evidence --- dies at the first
hash-ordered iteration.

\subsection{Worked scenario}\label{ch01:sec:nondet-scenario}

A well-meaning pipeline builds a carrier from a Python set literal:

\begin{quote}
\texttt{states = \{"calm", "alert", "overload"\}}\\
\texttt{carrier = tuple(states)}
\end{quote}

The programmer reasons: the set has three fixed elements, so the tuple has
three fixed elements; nothing here is random. The reasoning is wrong in one
load-bearing detail: \texttt{tuple(states)} fixes the \emph{membership} but
takes the \emph{order} from set iteration, and set iteration order for
strings is a function of the process's hash seed.

Run A of the pipeline, under one seed, obtains the carrier order
$(\texttt{alert}, \texttt{calm}, \texttt{overload})$: the element
\texttt{calm} has standard index $1$. An endomap $f$ intended to send every
state to \texttt{calm} is then constructed by looking up indices, producing
mapping tuple $(1, 1, 1)$. The pipeline serializes the object and arrow ---
carrier order and mapping tuple --- into a canonical byte string and records
its SHA-256 digest in an evidence log, alongside the conclusions drawn from
the run.

Run B, the auditor's replay on another day (a fresh process, hence a fresh
seed), obtains carrier order
$(\texttt{overload}, \texttt{alert}, \texttt{calm})$: now \texttt{calm} has
standard index $2$, the same intended map serializes with mapping tuple
$(2, 2, 2)$ and a permuted carrier sequence, and the recorded digest does not
match. The auditor is left with exactly the wrong dichotomy: either the
original evidence was corrupted, or the pipeline is nondeterministic --- and
no local inspection of the code distinguishes the benign cause (hash-order
iteration of a three-element set) from a malignant one (actual tampering or
an actual semantic bug). The evidentiary value of the digest is zero, not
reduced: a digest that honest replays fail to reproduce cannot certify
anything.

The failure generalizes beyond digests. If two subsystems each build ``the
same'' carrier from set iteration within one process, they may agree by
luck; across processes, machines, or language versions they will not.
Basis-dependent quantities --- the matrix of the lift $A_f$, the shortlex
rank of a word, the identity of the canonical representative selected by
normal-form computation --- all inherit the nondeterminism. A ``canonical''
form that varies by process is a contradiction in terms.

Chapter~\ref{ch03:chap} develops this scenario mechanically ---
what \texttt{PYTHONHASHSEED} actually controls, how randomization enters
\texttt{set} and \texttt{dict} iteration, why the problem is a general
software-evidence problem rather than a Python quirk --- and states the
kernel's ordered-input rule: every order-sensitive constructor input rejects
sets, mappings, and one-shot iterables outright, at construction time, with a
typed error. The caller must supply a replayable ordered sequence, which
makes the ordering decision explicit, reviewable, and versionable.

\subsection{The lesson}

Determinism is not an implementation nicety here; it is a precondition for
the entire evidence discipline of the program. The kernel's artifacts are
bound by SHA-256 hashes in a manifest, its test outcomes are recorded against
those hashes, and its formal seed audits exact source hashes before building.
All of that machinery presupposes that honest reruns reproduce bytes. A
single hash-ordered iteration anywhere upstream of a digest silently voids
the presupposition. The ordered-input rule is therefore enforced at the
lowest level of the kernel --- inside the constructors --- rather than in
documentation or review checklists, on the general principle that a property
this load-bearing must be impossible to violate accidentally.

\section{The scenarios replayed against the kernel}\label{ch01:sec:replay}

Before stating the design philosophy in the abstract, it is worth
replaying each scenario against the kernel's actual rules, because the
three outcomes have a common shape: the failing step does not execute
wrongly --- it does not execute at all, and the error it raises names the
exact judgment that failed.

\subsection{Scenario I under typed composition}

In the kernel, the summarization and reading steps are arrow values
carrying their full judgments:

\begin{quote}
\texttt{p : DriveProfile -> EvalSummary} \quad with mapping
$(0, 0, 2, 2, 2, 1, 3)$\\
\texttt{r : QuadState -> DriveProfile} \quad with mapping
$(4, 6, 5, 0)$
\end{quote}

The composite request \texttt{compose(r, p)} is checked before any index
is touched: the declared target object of \texttt{p} is
$(\texttt{EvalSummary}, (e_0, e_1, e_2, e_3))$, the declared source
object of \texttt{r} is $(\texttt{QuadState}, (q_0, q_1, q_2, q_3))$,
and the two are unequal because their names differ. The kernel raises
\texttt{CompositionTypeError} with a message of the form

\begin{quote}
\texttt{cannot compose p:DriveProfile->EvalSummary then}\\
\texttt{r:QuadState->DriveProfile}
\end{quote}

The number $5$ --- the smuggled image of $q_2$ that the untyped
implementation returned --- is never computed, because the guard runs
before the pointwise formula. Note also what the fix looks like on the
kernel's terms: if the analyst genuinely intends an identification of
evaluative labels with quadrant labels, she must write it as an explicit
arrow $\texttt{EvalSummary} \to \texttt{QuadState}$ and insert it into
the composite. The identification then exists as a reviewable object
with an author and a version --- exactly the audit record whose absence
made the original failure silent
(Remark~\ref{ch03:rem:cost} returns to this point).

\subsection{Scenario II under the endomap/path separation}

The alternating session is representable in the kernel in exactly one
way: as a path value,

\begin{quote}
\texttt{session = SemigroupoidPath((a, b, a, b, a, b))}
\end{quote}

which validates (all six arrows are endo on the same object, so all
adjacent pairs match) and offers \texttt{composite()} --- the honest
six-step composite, which evaluates to the transposition
$h^3 = h$ computed in Section~\ref{ch01:sec:power-scenario}. What the
path does \emph{not} offer is any power operation: there is no
\texttt{session.power(n)}, by design, and the runtime type
\texttt{SemigroupoidPath} is not accepted where an \texttt{Endomap} is
required. The analyst who wants to invoke the factorial idempotence
theorem must first construct \texttt{Endomap(g)} for some single arrow
\texttt{g} --- and for the alternating session no such single arrow
exists with the session's semantics, a fact the analyst now confronts at
the type level instead of discovering, twelve steps into production,
through a falsified stabilization claim. If she constructs
\texttt{Endomap(session.composite())} she gets a legitimate object ---
the one-round map $h$ composed out to six steps --- and every theorem she
then invokes is true \emph{of that map}, which is the point: the type
system has forced the quantifier back over single endomaps, where the
theorems actually hold.

The cross-carrier variant fails even earlier: a path whose adjacent
objects mismatch raises \texttt{CompositionTypeError} at construction,
naming the offending segment, and a mismatched \texttt{Endomap}
construction --- an arrow whose source and target are different objects
--- likewise raises at construction with a message of the form
\texttt{"g' has type X0->X2, so it is not an endomap"}.

\subsection{Scenario III under the ordered-input rule}

The pipeline's first line, transplanted into the kernel, is already
illegal:

\begin{quote}
\texttt{FiniteObject(name="States", elements=\{"calm", "alert",
"overload"\})}
\end{quote}

raises \texttt{ValidationError}: finite-object elements must be an
ordered sequence; sets, mappings, and one-shot iterables are not
accepted. The programmer is forced to write the order:

\begin{quote}
\texttt{FiniteObject(name="States", elements=("calm", "alert",
"overload"))}
\end{quote}

and from this point Invariant~\ref{ch03:inv:replay} takes over: the
constructed object, the arrow built over it, the serialization, and the
digest are pure functions of the written program text. Run A and Run B
now agree byte-for-byte, on any machine, under any
\texttt{PYTHONHASHSEED}, because the hash seed no longer participates in
any datum the kernel emits. The auditor's replay reproduces the digest,
and the digest is evidence again. The cost --- one explicit ordering
decision, visible in the diff --- is the entire price, and
Section~\ref{ch03:sec:hash-order} argues it is the minimum price any
solution can charge, since the ordering decision existed all along and
was merely being made, invisibly and differently each run, by the hash
table.

\subsection{The common shape}

In all three replays the kernel's contribution is the same: a judgment
that informal notation left implicit --- the middle objects of a
composite, the single-endomap hypothesis of an iteration theorem, the
enumeration order of a carrier --- is reified as explicit, checkable
data, and the check runs \emph{before} any computation that would depend
on the judgment being true. The failures do not become error codes to
handle; they become programs that cannot be written. This is what
``make the failure a type error'' means concretely, and the remainder of
the monograph is the systematic execution of that idea across the
kernel's five capabilities.

\section{Design philosophy of the kernel}\label{ch01:sec:philosophy}

Four commitments shape everything that follows.

\subsection{Smallest kernel, hardest boundary}

TKS v8.0 is the \emph{smallest} kernel that turns each of the three failures
into a type error. Smallness is a feature with teeth: the implementation is
three modules (\texttt{category.py}, \texttt{deterministic\_lift.py},
\texttt{normal\_forms.py}, totalling under 26 kilobytes of source), the
mathematics is elementary and completely proved, and the bounded exhaustive
test families can genuinely enumerate their domains ($19{,}683$
associativity triples, $289$ endomaps, $499$ typed maps, $1{,}678$
lift-composition pairs; Chapter~10). A larger kernel would dilute each of
these properties. In particular the kernel is an \emph{independent reference
model}: it does not modify, replace, or certify the active CAS notation
compiler, it has no production runtime dependency, and a manifest-boundary
regression test asserts exactly that.

\subsection{Refusal to encode canon before review}

The kernel deliberately encodes no Noetic, no Domain, no Drive, no Register,
and no CAS construct of any sort. This is the program's most frequently
misread decision, so we state its logic plainly. The mapping from the
canonical ontology onto a categorical substrate is itself a research claim
--- arguably the central one --- and admitting it by construction, as a side
effect of writing the kernel, would launder an unreviewed claim into
infrastructure. The mapping is instead held as an explicit open obligation
(PO-TKS8-CORE-08, \textsc{open / not licensed}), to be performed, if at all,
under separate review against the corrected v7.4 theorem records. Until
then, an arrow in this book means a total function between named finite
carriers, and nothing else.

\begin{boundary}\label{ch01:boundary:no-canon}
No result in this monograph is a claim about cognition, about the TKS/CAS
ontology, about the production compiler's semantics, or about any empirical
system. The kernel maps no canonical construct onto arrows. Where a scenario
in this chapter uses ontology-flavored names (\texttt{DriveProfile},
\texttt{EvalSummary}), the names are illustrative strings; the mathematics
sees only carriers and index tuples.
\end{boundary}

\subsection{Fail closed}

Wherever the kernel cannot deliver a complete, correct answer, it raises a
typed error rather than returning a best effort. Composition with mismatched
middle objects raises rather than coercing; constructors reject unordered
inputs rather than guessing an order; the normal-form closure raises
\texttt{ClosureLimitExceeded} rather than returning a partially explored
table as ``canonical.'' The alternative --- silently truncated canonical
forms, silently coerced compositions --- is precisely the raw material of the
three failure modes.

\subsection{Evidence as a first-class artifact}

Every mathematical claim in the program carries a registered evidential
status, and the statuses are heterogeneous on purpose: some claims have
complete informal proofs plus bounded exhaustive tests; some have partial
formal verification in a pinned Lean~4.32.1 seed \cite{lean4} whose verifier
admits the single standard axiom \texttt{propext} and rejects project axioms
and \texttt{sorry}; some are standard results cited to the literature with
formalization open; some are open outright. The full ledger is reproduced in
Appendix~B. The monograph never averages these statuses into an impression
of overall verifiedness: bounded testing is not proof-assistant verification
of the running code, the Lean seed's scope is exactly enumerated
(Chapter~11), and the equivalence between the Python implementation and the
Lean model is itself an open obligation. This stance follows the small-scope
tradition in software modeling \cite{jackson}: exhaustive checking of a
bounded model is a powerful falsifier and a weak verifier, and conflating the
two is how disciplines rot.

\section{Contributions and their evidential status}\label{ch01:sec:contributions}

The kernel supplies five capabilities. For each we state the capability, the
chapters that develop it, and the exact evidential status, quoted from the
ledger.

\begin{enumerate}
\item \emph{Named finite objects with explicitly ordered carriers}
(Chapter~\ref{ch03:chap}). Object identity is the pair of declared name and
ordered carrier; all order-sensitive constructor inputs reject sets,
mappings, and one-shot iterables. Status: executable validation ---
constructor and negative-fixture tests; this is a design discipline, not a
theorem.
\item \emph{Total finite-function arrows with partial, type-checked
composition satisfying the category laws} (Chapters~\ref{ch03:chap} and~4).
Status: partial formal verification present --- pointwise proof, bounded
exhaustive tests over $19{,}683$ associativity triples and $499$ typed maps,
and exact Lean lemmas for the Lean model; Python-to-Lean equivalence open.
\item \emph{Type separation of endomap powers from semigroupoid paths}
(Chapters~5 and~6), with the safe finite power theorem
$f^{|X|!} \circ f^{|X|!} = f^{|X|!}$. Status: the separation is executable
validation (distinct runtime types, negative tests); the theorem is a
standard finite-semigroup result \cite{pin, howie} verified exhaustively for
all $289$ endomaps on carriers through size four, with formalization open.
\item \emph{The deterministic standard-basis lift} $A_f e_x = e_{f(x)}$
(Chapter~7), with the exact norm formula
$\lVert A_f \rVert^2 = \max_{y} |f^{-1}(y)|$, its trace, rank, and isometry
corollaries, functoriality, and a permanent noncontractivity counterexample.
Status: informal proof present; bounded exhaustive structural tests
($1{,}678$ exact-matrix composition pairs); Lean kernel-checks the exact
two-state counterexample and the identity and basis-composition laws;
universal formalization and the finite-sum matrix bridge open.
\item \emph{Stable denotational normal forms} by shortlex minimization over a
completed finite closure (Chapter~8), with soundness of representative-first
pruning and fail-closed resource semantics. Status: informal correctness
argument present; regression tests; formalization open.
\end{enumerate}

Chapters~9 through~11 document the implementation, the executable evidence
(the validation snapshot of 2026-08-04 records \textsc{pass} on 20 tests),
and the formal seed. Chapter~12 states the integration boundaries and the
downstream program. Nothing in the list above is peer-reviewed
(PO-TKS8-CORE-10, \textsc{open}); no empirical, cognitive-dynamics,
channel-contractivity, or CAS-efficacy claim follows from any item; and the
mathematical content is standard --- the contribution claimed is the
\emph{engineering discipline} binding it together, not new theorems.

The correspondence between capabilities, obligations, and statuses is
summarized in Table~\ref{ch01:tab:status}; the authoritative record is
the ledger itself (Appendix~B), and where this table and the ledger could
ever disagree, the ledger governs.

\begin{table}[htbp]
\centering
\begin{tabular}{@{}llll@{}}
\toprule
Capability & Chapters & Ledger items & Status (abbreviated) \\
\midrule
Objects, ordered carriers & 3 & PO-02 (part) & executable validation \\
Typed composition, laws & 3, 4 & PO-01, PO-02 & partial formal verif. \\
Endomap/path separation & 5, 6 & PO-03, PO-04 & exec.\ valid.; std.\ result \\
Deterministic lift, norm & 7 & PO-05, PO-11 & informal proof; partial \\
Shortlex normal forms & 8 & PO-06 & informal argument \\
Lean formal seed & 11 & PO-07 & partial, present \\
Semantic mapping to CAS & 12 & PO-08 & open / not licensed \\
Compiler integration & 12 & PO-09 & open \\
Specialist review & --- & PO-10 & open \\
\bottomrule
\end{tabular}
\caption{Capabilities against ledger obligations (PO-TKS8-CORE numbers
abbreviated). ``Partial formal verif.''\ always means: exact Lean lemmas
about the Lean model, bounded exhaustive tests of the Python model, and
an open bridge between the two.}
\label{ch01:tab:status}
\end{table}

\section{Positioning among neighboring disciplines}\label{ch01:sec:positioning}

The kernel's ideas are borrowed, individually, from well-established
practice; the claimed contribution is their joint enforcement in one
small, evidence-bound artifact. Situating each borrow both credits the
sources and sharpens what the kernel does not attempt.

\emph{Typed composition} is the oldest idea here: that the legality of
combining two operations should be checked against declared interfaces
rather than discovered at runtime is the founding insight of static type
systems, and its categorical formulation --- composition as a partial
operation governed by objects --- is textbook category theory
\cite{maclanecwm}. The kernel adds only strictness choices: name-inclusive
object identity and the refusal of every coercion.

\emph{The endomap/path separation} is elementary semigroup theory
\cite{howie, pin} read as a typing discipline: powers live in
endomorphism monoids and nowhere else, so the type system should make
``nowhere else'' unrepresentable. The mathematical content
(Chapter~\ref{ch02:chap}) is classical; the engineering content is the
decision to encode a quantifier's hypothesis as a runtime type.

\emph{Determinism as an evidence requirement} is the reproducible-builds
insight applied to mathematical artifacts: a digest certifies nothing
unless honest replays reproduce it. The kernel's ordered-input rule is a
library-level instance of that doctrine, and
Chapter~\ref{ch03:chap} argues the general case.

\emph{Bounded exhaustive testing} follows the small-scope tradition
\cite{jackson}: enumerate \emph{all} instances of a small model rather
than sampling a large one, on the hypothesis that most defects manifest
at small scope. The kernel's test families are sized by exact counts
(Chapter~\ref{ch02:chap}, Section~\ref{ch02:sec:counting}) so that
``all'' is a theorem, not an aspiration.

\emph{Kernel-checked formalization} follows standard proof-assistant
practice \cite{lean4}: a small trusted checker, a pinned toolchain, an
audited axiom footprint. The program's honesty rule --- the Lean model
and the Python implementation are distinct artifacts whose equivalence
is an open obligation --- is itself standard practice in careful
verification work, and Chapter~11 maintains it without exception.

What the kernel does not attempt, even within these disciplines: it is
not a general-purpose proof assistant, not a rewriting engine, not a
model checker, and not a type-system extension of any existing language
--- it is a library whose constructors refuse, which is the smallest
mechanism that suffices for the three failures.

\subsection{Falsifiability of the kernel's claims}

A reference kernel should say what would refute it. Each capability has
a concrete falsifier: a composable triple of typed arrows on which the
implementation's composition disagrees with pointwise evaluation; a
constructible arrow accepted with mismatched middle objects; a
\texttt{SemigroupoidPath} value on which any power-like operation
succeeds; an endomap $f$ on a carrier of size at most four with
$f^{|X|!} \circ f^{|X|!} \neq f^{|X|!}$; a typed pair $(f, g)$ with
$A_{g \circ f} \neq A_g A_f$ as exact integer matrices; a catalog and
denotation for which the closure returns a word that is not
shortlex-least among well-typed representatives, or returns a partial
table without raising; and any kernel output that differs between two
runs with identical constructor arguments. The bounded exhaustive
families of Chapter~10 are precisely systematic searches for the first
five falsifiers within their bounds; none has been found, which is
evidence within the bounds and silence beyond them
(Remark~\ref{ch01:rem:status}).

\section{Reading guide}\label{ch01:sec:guide}

The monograph is organized so that a reader can enter at several depths.

\emph{Part I: Foundations.} Chapter~\ref{ch02:chap} is a self-contained
mathematical primer --- finite sets and total functions, counting, the
semigroup-to-category ladder, functional graphs of endomaps,
finite-dimensional Hilbert spaces, and shortlex orders --- cited throughout
as standard background \cite{bookofproof, maclanecwm, pin, howie, horn}. A
reader fluent in these topics may skim it for notation.
Chapter~\ref{ch03:chap} defines objects and arrows and develops the
determinism discipline at full depth.

\emph{Part II: The kernel's mathematics.} Chapter~4 proves the category
laws. Chapter~5 develops endomap theory and the safe finite power.
Chapter~6 develops semigroupoid paths and the impossibility of path powers.
Chapter~7 develops the lift: norm theorem, corollaries, functoriality, and
the permanent counterexample. Chapter~8 develops normal forms: shortlex
minimization, the representative-first closure algorithm, its soundness, and
its resource semantics.

\emph{Part III: Evidence.} Chapter~9 walks the implementation. Chapter~10
specifies the bounded exhaustive test families and what they do and do not
establish. Chapter~11 documents the pinned Lean seed, its axiom audit, and
its exact verified and open scope.

\emph{Part IV: Boundaries.} Chapter~12 states what is excluded, the three
downstream boundary objects that consume the kernel strictly through its
types, and the reproduction procedure.

\emph{Appendices.} Appendix~A collects worked examples end to end.
Appendix~B reproduces the complete proof-obligation ledger
(PO-TKS8-CORE-01 through 11) with statuses. Appendix~C is a glossary.

Readers primarily interested in \emph{why} the kernel exists should read this
chapter and Chapter~12. Readers auditing the mathematics should read
Chapters~\ref{ch02:chap} through~8. Readers replaying the evidence should
read Chapters~9 through~11 and Appendix~B.

\section{Notation and conventions}\label{ch01:sec:notation}

\begin{notation}\label{ch01:not:conventions}
Throughout the monograph:
\begin{itemize}
\item Finite objects are written $X, Y, Z$; $|X|$ is the carrier length;
carrier positions are the \emph{standard indices} $0, \dots, |X|-1$.
\item Arrows are written $f : X \to Y$; $\id_X$ is the identity;
$g \circ f$ is composition, defined iff the target object of $f$ equals the
source object of $g$; $\im(f)$ is the image; $f^{-1}(y)$ is the fiber over
$y$.
\item $\mathrm{End}(X)$ is the monoid of endomaps of $X$; $f^n$ is the
$n$-th monoid power, $f^0 = \id_X$.
\item $H_X$ is the complex Hilbert space with orthonormal basis
$(e_x)_{x \in X}$ indexed by the ordered carrier; $A_f$ is the deterministic
standard-basis lift of $f$; $A_f^{*}$ its adjoint; $\lVert \cdot \rVert$
the operator norm; $\trace$ and $\rk$ the trace and rank.
\item $\llbracket w \rrbracket$ is the denotation of a typed word $w$;
$\NF_G(a)$ is the shortlex-least well-typed word denoting $a$ relative to
the declared generator catalog $G$; $<_{\mathrm{slex}}$ is the shortlex
order.
\item Code identifiers appear in typewriter face
(\texttt{FiniteObject}, \texttt{denotation\_key}); they name entities of the
reference implementation in \texttt{src/tks\_v8\_core}.
\item Claim-status words --- \textsc{open}, \textsc{not licensed},
\textsc{partial formal verification present}, \textsc{executable
validation}, \textsc{standard result} --- are quoted from the obligation
ledger and are never paraphrased.
\end{itemize}
\end{notation}

Table~\ref{ch01:tab:symbols} collects the recurring symbols with the
location of their governing definitions.

\begin{table}[htbp]
\centering
\begin{tabular}{@{}lll@{}}
\toprule
Symbol & Meaning & Defined \\
\midrule
$X = (\mathrm{name}, \mathrm{carrier})$ & finite object & Def.~\ref{ch03:def:object} \\
$|X|$ & carrier length & Def.~\ref{ch03:def:object} \\
$f : X \to Y$ & arrow judgment & Def.~\ref{ch03:def:arrow} \\
$\id_X$, $g \circ f$ & identity, composition & Ch.~2, Ch.~4 \\
$\im(f)$, $f^{-1}(y)$ & image, fiber & Def.~\ref{ch02:def:image-fiber} \\
$\mathrm{End}(X)$, $f^{n}$ & endomorphism monoid, power & Prop.~\ref{ch02:prop:ladder} \\
$T_X$ & full transformation monoid & Def.~\ref{ch02:def:full-transf} \\
$H_X$, $(e_x)_{x \in X}$ & lift space, standard ONB & Ch.~2 Sec.~7, Ch.~7 \\
$A_f$, $A_f^{*}$ & lift, adjoint & Ch.~7 \\
$\lVert \cdot \rVert$, $\trace$, $\rk$ & operator norm, trace, rank & Def.~\ref{ch02:def:opnorm} \\
$J_m$ & all-ones $m \times m$ matrix & Prop.~\ref{ch02:prop:allones} \\
$\llbracket w \rrbracket$ & denotation of a typed word & Ch.~8 \\
$<_{\mathrm{slex}}$, $\NF_G$ & shortlex order, normal form & Def.~\ref{ch02:def:shortlex}, Ch.~8 \\
$\rk_{\mathrm{slex}}(w)$ & shortlex rank of a word & Prop.~\ref{ch02:prop:slex-rank} \\
$\cC$, $\Ob$, $\Hom$ & category, objects, hom-sets & Def.~\ref{ch02:def:category} \\
\bottomrule
\end{tabular}
\caption{Recurring symbols. References into Chapters 4--8 point at the
governing definitions of Part II; primer references are exact.}
\label{ch01:tab:symbols}
\end{table}

\begin{remark}[How to read the status vocabulary]\label{ch01:rem:status}
\textsc{Executable validation} means: the property is enforced by
constructors and demonstrated by tests, including negative fixtures; it is a
statement about the finite executable model only. \textsc{Standard result}
means: the mathematics is textbook material, cited as background and proved
here for self-containment, with no novelty claimed. \textsc{Partial formal
verification present} means: a pinned proof assistant kernel-checks stated
lemmas about a Lean model of the definitions, under an audited axiom
footprint, and the bridge from that model to the running Python is open.
\textsc{Open} means open: no evidence is claimed. \textsc{Not licensed}
means that even drafting an argument for the claim is out of scope for this
lane until a separate review admits it.
\end{remark}

With the failures, the philosophy, and the notation in hand, we turn to the
mathematics. Everything in the next chapter is standard; everything after it
is discipline.
